\begin{frame}[t]{Programar começa com uma regra clara}
Uma estufa usa um sinalizador para indicar temperatura elevada.
\begin{itemize}
\item Entrada: temperatura medida em graus Celsius.
\item Regra: sinalizar quando a temperatura for \textbf{maior ou igual a 30}.
\item Saída: sinalizador ligado ou desligado.
\end{itemize}
\begin{alertblock}{Antes de escrever código}
``Ligar quando estiver quente'' é ambíguo. O valor limite e o comportamento exatamente no limite precisam estar definidos.
\end{alertblock}
\end{frame}
